\section{The exact gapless auxiliary operator and the physical boundary}
\label{cusp-interpretation-section}

The preceding flat-torus calculation produces eigenvalues tending to zero
along the radial cusp.  This section identifies a single self-adjoint
operator for which those numbers genuinely form a gapless spectrum, proves
that complex and Hodge-theoretic data alone cannot determine such a
spectrum, and states the exact additional morphism needed to transfer an
auxiliary Rayleigh sequence to a physical Hamiltonian.  In particular, the
word \emph{gapless} below is attached to a specified operator and Hilbert
space, not to the complex threefold in isolation.

\subsection{A direct-integral scalar Laplacian with zero essential threshold}
\label{cusp-interpretation-direct-integral-subsection}

Use the angular coordinates \(y\in\mathbb R^4/\mathbb Z^4\) of
\cref{cusp-spec-exact-dual-lattice}.  The map \(y\mapsto P_Ty\) identifies
this fixed torus with \(F_T\), and normalized Haar measure on \(F_T\) pulls
back to \(dy\).  Write
\[
 e_k(y):=e^{2\pi i k\cdot y},
 \qquad
 \lambda_k(T):=4\pi^2\lvert P_T^{-\mathsf T}k\rvert^2,
 \qquad k\in\mathbb Z^4.
\]
By \cref{cusp-spec-exact-dual-lattice}, these are exactly the
scalar-Laplacian eigenfunctions and eigenvalues on \((F_T,g_T)\), normalized
in \(L^2\) for Haar probability measure.

Fix \(T_0\) large enough that all radial formulas of
\cref{cusp-spec-exact-dual-lattice,cusp-spec-shortest-modes} hold, and put
\begin{equation}
 \mathscr H_{\mathrm{aux}}
 :=L^2\!\left([T_0,\infty),dT;
 L^2_0(\mathbb R^4/\mathbb Z^4,dy)\right).
 \label{cusp-interpretation-auxiliary-Hilbert-space}
\end{equation}
Fourier expansion gives the concrete unitary identification
\begin{equation}
 \mathscr H_{\mathrm{aux}}
 \cong
 \bigoplus_{k\in\mathbb Z^4\setminus\{0\}}
 L^2([T_0,\infty),dT).
 \label{cusp-interpretation-Fourier-direct-sum}
\end{equation}
Define the multiplication operator
\begin{equation}
 (\mathscr L_{\mathrm{aux}}u)_k(T)
 :=\lambda_k(T)u_k(T)
 \label{cusp-interpretation-direct-integral-operator}
\end{equation}
on the exact domain
\begin{equation}
 D(\mathscr L_{\mathrm{aux}})
 :=\left\{u\in\mathscr H_{\mathrm{aux}}:
 \sum_{k\ne0}\int_{T_0}^{\infty}
 \lambda_k(T)^2\lvert u_k(T)\rvert^2\,dT<\infty\right\}.
 \label{cusp-interpretation-direct-integral-domain}
\end{equation}
Thus \(\mathscr L_{\mathrm{aux}}\) is, by definition and not by analogy,
the direct integral of the mean-zero scalar Laplacians \(\Delta_{F_T}\).

\begin{theorem}[A proved gapless auxiliary operator]
\label{cusp-interpretation-direct-integral-gapless}
The operator \(\mathscr L_{\mathrm{aux}}\) in
\cref{cusp-interpretation-direct-integral-operator} is nonnegative and
self-adjoint.  It has no zero eigenvector, but
\begin{equation}
 0\in\sigma_{\mathrm{ess}}(\mathscr L_{\mathrm{aux}}).
 \label{cusp-interpretation-zero-essential-spectrum}
\end{equation}
In particular, this specified auxiliary operator has no positive spectral
gap above zero.
\end{theorem}

\begin{proof}
For each \(k\ne0\), \(\lambda_k(T)\) is a positive continuous real function.
Multiplication by it on
\[
 \{v\in L^2([T_0,\infty)):\lambda_kv\in L^2([T_0,\infty))\}
\]
is self-adjoint and nonnegative: its adjoint domain is obtained by testing
against compactly supported bounded functions, which forces
\(\lambda_kv\in L^2\), and the adjoint acts by the same real multiplier.
The orthogonal sum of these multiplication operators has precisely the
domain in \cref{cusp-interpretation-direct-integral-domain}; hence it is
self-adjoint and nonnegative.  Since every \(\lambda_k(T)\) is strictly
positive, \(\mathscr L_{\mathrm{aux}}u=0\) implies \(u_k=0\) almost
everywhere for every \(k\ne0\), so its kernel is zero.

Let \(k_0=(1,0,0,0)\).  The exact formula
\cref{cusp-spec-unit-mode-exact-lengths} gives
\begin{equation}
 \lambda_{k_0}(T)
 =4\pi^2\frac{L_T^2+m_T^2}{D_T^2}
 =\frac{4\pi^2}{T^2}\bigl(1+O(T^{-1})\bigr).
 \label{cusp-interpretation-soft-mode-eigenvalue}
\end{equation}
Choose integers \(n>T_0\), and define \(u^{(n)}\) by
\begin{equation}
 u^{(n)}_{k_0}(T)=\mathbf1_{[n,n+1]}(T),
 \qquad u^{(n)}_k(T)=0\quad(k\ne k_0).
 \label{cusp-interpretation-Weyl-sequence}
\end{equation}
The vectors \(u^{(n)}\) lie in
\(D(\mathscr L_{\mathrm{aux}})\), are orthonormal, and converge weakly to
zero because their \(T\)-supports escape to infinity.  Moreover,
\[
 \lVert\mathscr L_{\mathrm{aux}}u^{(n)}\rVert^2
 =\int_n^{n+1}\lambda_{k_0}(T)^2\,dT
 \le\sup_{T\in[n,n+1]}\lambda_{k_0}(T)^2
 \longrightarrow0
\]
by \cref{cusp-interpretation-soft-mode-eigenvalue}.  This orthonormal
singular sequence proves
\(0\in\sigma_{\mathrm{ess}}(\mathscr L_{\mathrm{aux}})\).  Explicitly, if
zero were outside the essential spectrum, then, because the kernel is zero,
some interval \([0,\varepsilon)\) would have finite-dimensional spectral
projection.  Its complement would satisfy
\(\lVert\mathscr L_{\mathrm{aux}}v\rVert\ge
\varepsilon\lVert v\rVert\); weak convergence would make the components of
\(u^{(n)}\) in the finite-dimensional spectral subspace tend strongly to
zero, contradicting
\(\lVert\mathscr L_{\mathrm{aux}}u^{(n)}\rVert\to0\) and
\(\lVert u^{(n)}\rVert=1\).
\end{proof}

\begin{remark}[The unit-volume normalization is a different operator family]
The theorem uses the unscaled Euclidean quotient metrics \(g_T\).  By
\cref{cusp-spec-unit-volume}, the first positive scalar eigenvalue for
\(\widehat g_T=D_T^{-1/2}g_T\) still tends to zero, but at the different
rate
\[
 \lambda_1(F_T,\widehat g_T)
 =\frac{4\pi^2}{T}\bigl(1+O(T^{-1})\bigr),
\]
rather than the unscaled \(T^{-2}\) rate.  Rescaling a metric changes its
Laplacian, so the two operator families and their asymptotic scales must not
be identified.
\end{remark}

\subsection{Complex data do not select a spectral gap}
\label{cusp-interpretation-metric-indeterminacy-subsection}

\begin{proposition}[Fixed complex torus, fixed volume, arbitrarily small gap]
\label{cusp-interpretation-fixed-complex-counterexample}
Let
\[
 A=(\mathbb C/(\mathbb Z+i\mathbb Z))^2
\]
with fixed complex coordinates
\(z_j=x_j+iy_j\).  For \(R\ge1\), the metric
\begin{equation}
 g_R=R^2(dx_1^2+dy_1^2)+R^{-2}(dx_2^2+dy_2^2)
 \label{cusp-interpretation-fixed-complex-metric}
\end{equation}
is a flat K\"ahler metric of volume one on the same complex torus, while
\begin{equation}
 \lambda_1(\Delta_{g_R})=\frac{4\pi^2}{R^2}\longrightarrow0.
 \label{cusp-interpretation-fixed-complex-gap}
\end{equation}
Consequently the complex structure, Hodge decomposition, integral
cohomology, and coherent direct-image data of a trivial family do not, even
together with total volume, determine a positive lower bound for the scalar
Laplacian gap.
\end{proposition}

\begin{proof}
The associated two-form
\[
 \omega_R=R^2dx_1\wedge dy_1+R^{-2}dx_2\wedge dy_2
\]
is closed, positive, and \(J\)-invariant, so \(g_R\) is K\"ahler for the fixed
complex structure.  In the ordered real coordinates
\((x_1,y_1,x_2,y_2)\), its matrix is
\(\operatorname{diag}(R^2,R^2,R^{-2},R^{-2})\), whose determinant is one;
therefore the volume is one.  The Fourier character
\[
 e_{a,b,c,d}
 =\exp(2\pi i(ax_1+by_1+cx_2+dy_2))
\]
has eigenvalue
\[
 4\pi^2\left(R^{-2}(a^2+b^2)+R^2(c^2+d^2)\right).
\]
For \(R\ge1\), the minimum over nonzero integer quadruples is
\(4\pi^2R^{-2}\), attained at the four vectors
\[
 (a,b,c,d)\in
 \{(1,0,0,0),(-1,0,0,0),(0,1,0,0),(0,-1,0,0)\}.
\]
This proves
\cref{cusp-interpretation-fixed-complex-gap}.  None of the listed complex or
cohomological objects depends on \(R\), and for a trivial holomorphic family
their coherent direct images are likewise unchanged.
\end{proof}

\subsection{The exact bridge to a physical Hamiltonian}
\label{cusp-interpretation-physical-bridge-subsection}

The mathematical mass gap in the Clay formulation is a property of the
Hamiltonian of a constructed four-dimensional quantum Yang--Mills theory,
not of an internal scalar Laplacian or of a classical Hessian
\cite{JaffeWitten2000}.  Euclidean input reaches such a Hamiltonian only
after a positive-metric reconstruction satisfying, in particular, the
Osterwalder--Schrader requirements
\cite{OsterwalderSchrader1973,OsterwalderSchrader1975}.  The following
proposition isolates the exact operator-theoretic statement that any proposed
transfer must prove.

\begin{proposition}[Typed Rayleigh-sequence transfer]
\label{cusp-interpretation-typed-physical-transfer}
Let \(H\ge0\) be self-adjoint on a Hilbert space
\(\mathcal H_{\mathrm{phys}}\).  Define
\begin{equation}
 \mathcal V_0:=\ker H,
 \qquad P_0:=\mathbf1_{\{0\}}(H),
 \qquad
 \mathcal H_{\mathrm{exc}}:=(I-P_0)\mathcal H_{\mathrm{phys}},
 \label{cusp-interpretation-excitation-space}
\end{equation}
and let
\begin{equation}
 \mathfrak h[\psi]:=\lVert H^{1/2}\psi\rVert^2,
 \qquad
 \mathcal Q(H):=D(H^{1/2}).
 \label{cusp-interpretation-Hamiltonian-form}
\end{equation}
For each \(n\), let \(\mathcal K_n\) be a Hilbert space,
\(\mathfrak q_n\) a nonnegative quadratic form on a domain
\(D(\mathfrak q_n)\subseteq\mathcal K_n\), and
\(0\ne u_n\in D(\mathfrak q_n)\).  Suppose there is a linear map
\[
 J_n:\operatorname{span}\{u_n\}\longrightarrow\mathcal Q(H),
\]
and put
\(\psi_n=(I-P_0)J_nu_n\).  If there are \(c>0\), \(C_n\ge0\), and
\(\varepsilon_n\ge0\) such that
\begin{align}
 \lVert\psi_n\rVert
 &\ge c\lVert u_n\rVert,
 \label{cusp-interpretation-norm-noncollapse}\\
 \mathfrak h[\psi_n]
 &\le C_n\mathfrak q_n[u_n]
       +\varepsilon_n\lVert u_n\rVert^2,
 \label{cusp-interpretation-energy-transfer}\\
 C_n\frac{\mathfrak q_n[u_n]}{\lVert u_n\rVert^2}
   +\varepsilon_n&\longrightarrow0.
 \label{cusp-interpretation-physical-smallness}
\end{align}
Then the excitation Hamiltonian
\[
 H_{\mathrm{exc}}
 :=H\big|_{D(H)\cap\mathcal H_{\mathrm{exc}}}
\]
is self-adjoint and satisfies
\begin{equation}
 \inf\sigma(H_{\mathrm{exc}})=0,
 \qquad
 0\in\sigma_{\mathrm{ess}}(H_{\mathrm{exc}}).
 \label{cusp-interpretation-physical-gapless-conclusion}
\end{equation}
\end{proposition}

\begin{proof}
The spectral projection \(P_0\) commutes with \(H\) and \(H^{1/2}\), so
\[
 \mathcal H_{\mathrm{phys}}
 =\ker H\oplus\mathcal H_{\mathrm{exc}}
\]
is a reducing decomposition.  Hence \(H_{\mathrm{exc}}\) is self-adjoint,
nonnegative, has form domain
\(\mathcal Q(H)\cap\mathcal H_{\mathrm{exc}}\), and has zero kernel.
Define \(w_n=\psi_n/\lVert\psi_n\rVert\).  Then
\(\lVert w_n\rVert=1\),
\(w_n\in\mathcal Q(H)\cap\mathcal H_{\mathrm{exc}}\), and
\begin{align*}
 \mathfrak h[w_n]
 &=\frac{\mathfrak h[\psi_n]}{\lVert\psi_n\rVert^2}\\
 &\le c^{-2}\left(
 C_n\frac{\mathfrak q_n[u_n]}{\lVert u_n\rVert^2}
 +\varepsilon_n\right)\longrightarrow0.
\end{align*}
The variational characterization of the bottom of a nonnegative
self-adjoint operator gives
\[
 \inf\sigma(H_{\mathrm{exc}})
 =\inf_{0\ne\psi\in
 \mathcal Q(H)\cap\mathcal H_{\mathrm{exc}}}
 \frac{\mathfrak h[\psi]}{\lVert\psi\rVert^2}=0.
\]
Since the spectrum is closed, zero belongs to it; since
\(\ker H_{\mathrm{exc}}=0\), zero is not an eigenvalue.  It is therefore a
continuous spectral point, and hence belongs to the essential spectrum under
the standard definition that removes isolated eigenvalues of finite
multiplicity.
\end{proof}

\begin{remark}[Why one vacuum vector is not enough]
If the vacuum is unique, choose a normalized vector
\(\Omega\) spanning \(\ker H\).  Then
\(P_0=|\Omega\rangle\langle\Omega|\).  Without uniqueness, merely requiring
\(J_nu_n\perp\Omega\) does not project away every zero-energy state.  For
example,
\[
 H=\operatorname{diag}(0,0,1),\qquad
 \Omega=(1,0,0),\qquad \psi=(0,1,0)
\]
has \(\psi\perp\Omega\) and zero Rayleigh quotient, although the genuine
excitation space \((\ker H)^\perp\) is the third coordinate and has gap one.
Thus \cref{cusp-interpretation-excitation-space} must use the full kernel
unless vacuum uniqueness is independently proved.
\end{remark}

\begin{auditrecord}[What the \(S^6\) calculation proves]
\label{cusp-interpretation-present-boundary}
The radial period matrix, scalar spectra, classical gauge Hessian, finite
cochain operator, magnetic Laplacian, and reducible \(SU(2)\) Jacobi operator
are explicit auxiliary geometric or classical operators.  They prove the
operator statements made in
\cref{cusp-interpretation-direct-integral-gapless,cusp-gauge-continuum-gap,cusp-gauge-finite-gap-limits,cusp-magnetic-oscillator-theorem,cusp-magnetic-negative-mode}.
They do not construct a common gauge-invariant positive-metric physical
Hilbert space \(\mathcal H_{\mathrm{phys}}\), a quantum Hamiltonian \(H\),
or maps \(J_n\) satisfying
\cref{cusp-interpretation-norm-noncollapse,cusp-interpretation-energy-transfer,cusp-interpretation-physical-smallness}
after continuum, infinite-volume, and renormalization limits.  Nor do the
period, Hodge, Leray, or direct-image calculations determine those data, as
\cref{cusp-interpretation-fixed-complex-counterexample} already disproves at
the level of geometric Laplacians.  Therefore the auxiliary gaplessness is
\Verified, while a conclusion about the four-dimensional quantum
Yang--Mills Hamiltonian is \Open.
\end{auditrecord}
